\chapter{Introduction}
\label{ch:intro}

An extra robot arm worn on the body is not a prosthesis. The person keeps the
limbs they were born with, and the machine works alongside
them~\cite{yang2021srlreview,prattichizzo2021augmentation}. The idea is old
enough to have a literature and young enough that most of it is still about
building the hardware rather than living with it.

What makes these devices awkward to control is easy to state. A robot arm
bolted to a bench has a base that stays where it was put. A robot arm bolted
to a person has a base that breathes, sways, turns to look at something and
occasionally flinches. It is also soft, irreplaceable, and within arm's reach
of the thing being moved.

This project works on the platform at Imperial College London known as Dr.\
Octopus~\cite{muve}, which carries two Kinova Gen3 seven-jointed
arms~\cite{kinova2022gen3} on a backpack frame. The arms are worn by one
person and driven by another. That arrangement is the starting point rather
than a detail, and \cref{fig:overview} shows how the parts sit together.

\begin{figure}[htbp]
  \centering
  % Who is where, and what passes between them.
\begin{tikzpicture}[x=1mm, y=1mm, font=\scriptsize]

  \node[person, minimum height=22mm, text width=30mm, fill=violet!10] (drv)
        at (0,0) {\textbf{The driver}\\[2pt] stands across the room, moves a
        sensed mannequin arm, or holds two hand controllers, or types a
        sentence};

  \node[stage, minimum height=22mm, text width=32mm] (sw) at (48,0)
        {\textbf{The software}\\[2pt] turns the input into a wanted hand
        position, checks it against the wearer's body, and sends joint
        commands};

  \node[person, minimum height=22mm, text width=30mm, fill=violet!10] (wr)
        at (96,0) {\textbf{The wearer}\\[2pt] carries \SI{17}{\kilo\gram} on a
        backpack frame and does not control the arms};

  \node[motor, minimum height=10mm, text width=32mm] (arms) at (48,-26)
        {\textbf{Two seven-jointed arms}\\ mounted behind the wearer's
        shoulders};

  \draw[go] (drv) -- node[tag, above]{intent} (sw);
  \draw[go] (sw) -- (arms);
  \draw[go] (arms.east) -- ++(20,0) |- node[tag, right, pos=0.25]
        {motion beside the head} (wr.south);
  \draw[go, violet!60!black] (wr.north) -- ++(0,10) -| node[tag, above, pos=0.25]
        {where the body is} (sw.north);
  \draw[go, dashed] (arms.west) -- ++(-22,0) |- node[tag, left, pos=0.25]
        {what the cameras see} (drv.south);

  \node[guard, minimum height=8mm, text width=44mm] (stop) at (48,-42)
        {\textbf{A third person} holds the emergency stop and watches the
        wearer, not the screen};
  \draw[halt] (stop) -- (arms);

\end{tikzpicture}

  \caption[Who is where]{The arrangement this thesis is about. Two people, not
  one. The driver commands the arms and cannot feel them; the wearer carries
  them and does not command them. The software between the two is responsible
  for the wearer's safety, because neither person is in a position to
  guarantee it.}
  \label{fig:overview}
\end{figure}

\section{The problem}

Two arms of seven joints each give fourteen things that can move. A person has
two hands and a limited amount of attention. Closing that gap is the design
problem, and there is no way to do it that does not give something up.

The obvious approach is to give the driver a small copy of the robot and let
them pose it. That is the approach taken here, and it works up to a point.
The mannequin master reproduces the sequence of joints in the real arm at
roughly a third of the size. What it cannot reproduce is the accuracy: it
routes fourteen analogue signals through joints that turn, and wires that
flex through turning joints eventually stop reporting. By the time this
project was measuring rather than building, seven of the fourteen channels no
longer produced a usable signal.

That could be treated as a repair job. It was instead treated as the normal
operating condition, because the same thing will happen again to whichever
channels are working next month. The system therefore has to answer three
questions at all times: which channels can be believed, what the arm can still
be commanded to do with the ones that remain, and what it should do when the
answer is nothing.

The second problem is the wearer. Standard robot software checks whether a
planned motion would touch anything, and reports either contact or no contact.
A person standing inside a pair of moving arms needs more than that. They need
a margin, held all the way along the path and not only at the ends, and they
need it enforced against the parts of the arm that actually swing past them,
which are the ones nearest the shoulder rather than the hand. Getting this
right turned out to be harder than expected, and \cref{ch:safety} is largely
an account of the ways it can be got wrong while appearing to work.

\section{Aims}

\begin{enumerate}
  \item Build a sensed mannequin arm that reproduces the robot's shape at
    reduced size, and characterise what its sensors actually deliver rather
    than what they are specified to deliver.
  \item Turn the driver's hand movement into arm commands in a way that keeps
    working as sensor channels fail, and measure what each failure costs.
  \item Build a safety layer suited to a machine mounted beside a person, and
    test it by deliberately breaking things rather than by reading the code.
  \item Measure where the arms can actually work once the wearer's body is
    taken into account, and say plainly which tasks that rules out.
  \item Add a second and third way of driving the arms, so that the study has
    something to compare against, and check that each reaches the same
    commands.
  \item Design the two-person study in full, and identify by measurement what
    has to be fixed before it can be run.
\end{enumerate}

\section{What is new here}

Three things in this thesis are not in the work it builds on.

The first is the treatment of sensor loss as a designed-for condition. Systems
of this kind are usually described in their working state. Here each channel
carries a health verdict computed from its own recent behaviour, the checks
applied to an incoming frame depend on which channels the current mode
actually reads, any channel can be switched off while running, and the cost of
switching it off is published rather than assumed. When too little is left, the
system reports that it has no position rather than inventing one. On a machine
beside somebody's head, an invented position is worse than an admitted gap.

The second is a negative result with a price attached. Under the mount as
built, the two arms share no reachable point anywhere in front of the wearer.
That is not a tuning failure and no amount of better control will change it. It
is reported here with the four independent measurements that establish it and
with a costed mechanical change that would improve it.

The third is methodological, and it is the part most likely to be useful
elsewhere. Seventeen times during this project, a result that looked like a
fault in the robot turned out to be a fault in the thing measuring the robot.
\Cref{ch:results} reports the numbers; the habit that produced them, and the
cases that taught it, are set out in \cref{sec:instruments}.

\section{How this thesis is arranged}

\Cref{ch:background} covers the work this builds on and states what is left
open. \Cref{ch:system} describes the system as built, in overview.
\Cref{ch:master} covers the mannequin arm and how a hand movement becomes an
arm command. \Cref{ch:safety} covers the safety layer.
\Cref{ch:perception} covers the cameras, and the change from a robot that was
told where things are to one that measures where they are.
\Cref{ch:study} sets out the study: tasks, conditions, measures and ethics.
\Cref{ch:results} reports everything measured on the system itself.
\Cref{ch:participant} is where the results from people will go, and is empty.
\Cref{ch:discussion} weighs it up, and \cref{ch:conclusion} says what should
happen next and in what order.
